\section{Comparing localization across contexts}
\label{sec:context-comparison}

Support restriction gives an exact comparison when one regime is contained in another. Cross-context studies are often harder: two prompt classes or intervention regimes may overlap only partially, or not at all at the input level. The relevant overlap for a mechanistic claim can nevertheless occur later in the network.

\subsection{Input overlap can understate internal overlap}

Fix an ambient source regime $C_\ast$. For $C,D\subseteq C_\ast$, the internal supports at cut $q$ are
\[
S_q^{\theta,C}=A_q^\theta(C),
\qquad
S_q^{\theta,D}=A_q^\theta(D).
\]
Their shared source inputs generate only
\[
S_q^{\theta,C\cap D}=A_q^\theta(C\cap D).
\]
There can be additional shared internal states because the prefix $A_q^\theta$ need not be injective.

\begin{proposition}[Source intersection and internal overlap]
\label{prop:source-overlap}
For $C,D\subseteq C_\ast$ and every cut $q$,
\begin{equation}
S_q^{\theta,C\cap D}
\subseteq
S_q^{\theta,C}\cap S_q^{\theta,D}.
\label{eq:source-overlap-inclusion}
\end{equation}
If $A_q^\theta|_{C_\ast}$ is injective, equality holds for every $C,D\subseteq C_\ast$. If it is not injective, there exist $C,D\subseteq C_\ast$ for which the inclusion is strict.
\end{proposition}

\begin{proof}
The inclusion is the general relation $f(C\cap D)\subseteq f(C)\cap f(D)$ applied to $f=A_q^\theta$. If $A_q^\theta|_{C_\ast}$ is injective and $z$ lies in both images, then $z=A_q^\theta(x)=A_q^\theta(x')$ for some $x\in C$ and $x'\in D$. Injectivity gives $x=x'\in C\cap D$, so $z\in S_q^{\theta,C\cap D}$.

If the restriction is noninjective, choose distinct $x,x'\in C_\ast$ with $A_q^\theta(x)=A_q^\theta(x')=z$. Taking $C=\{x\}$ and $D=\{x'\}$ gives $C\cap D=\varnothing$ while $z\in S_q^{\theta,C}\cap S_q^{\theta,D}$.
\end{proof}

Thus two disjoint prompt classes can still pass through a common internal state. Comparing explanations only on the image of their shared input labels can miss part of the domain on which the internal claims genuinely overlap.

\subsection{Cross-support compatibility and union-level realization}

A localization can work on two regimes separately and fail when the regimes are combined. The obstruction can be stated exactly. Let $S_1,S_2\subseteq X$ be analysis supports, let
\[
L:S_1\cup S_2\twoheadrightarrow U
\]
be a fixed access corestricted to its realized image on the union, and let
\[
\Phi:S_1\cup S_2\to Y_\Phi
\]
be a fixed phenomenon. For $i\in\{1,2\}$, corestrict $L|_{S_i}$ to $L(S_i)$ when evaluating realization on $S_i$.

\begin{proposition}[Cross-support compatibility criterion]
\label{prop:cross-support-compatibility}
Assume $L|_{S_1}$ realizes $\Phi|_{S_1}$ and $L|_{S_2}$ realizes $\Phi|_{S_2}$. Then $L$ realizes $\Phi$ on $S_1\cup S_2$ if and only if
\begin{equation}
x\in S_1,\ y\in S_2,\ L(x)=L(y)
\quad\Longrightarrow\quad
\Phi(x)=\Phi(y).
\label{eq:cross-support-compatibility}
\end{equation}
\end{proposition}

\begin{proof}
By local realization, every pair in $S_1\times S_1$ or $S_2\times S_2$ that lies in $\Ker L$ already lies in $\Ker\Phi$. The only access-fiber pairs on $S_1\cup S_2$ not covered by those two local conditions are cross-support pairs. Hence
\[
\Ker(L|_{S_1\cup S_2})
\subseteq
\Ker(\Phi|_{S_1\cup S_2})
\]
holds exactly when every $x\in S_1$ and $y\in S_2$ with $L(x)=L(y)$ also satisfy $\Phi(x)=\Phi(y)$. The claim follows from \cref{thm:realization-criterion}.
\end{proof}

The condition checks the new comparisons introduced by the union. Local realization handles pairs internal to each support; union-level realization adds the requirement that access collisions across the two supports remain phenomenon-consistent.

\begin{example}[Local realization need not extend to a union]
\label{ex:local-not-global}
Let
\[
Q(a,b)=a,
\qquad
S_1=\{(0,0),(1,1)\},
\qquad
S_2=\{(0,0),(2,1)\},
\]
and take the access $L=\pi_{\{2\}}$. On $S_1$, the second coordinate determines $Q$: $b=0$ gives $Q=0$ and $b=1$ gives $Q=1$. On $S_2$, it again determines $Q$: $b=0$ gives $Q=0$ and $b=1$ gives $Q=2$. Thus $L$ realizes $\Phi=Q$ on each support, and trivially on their intersection.

On the union,
\[
S_1\cup S_2
=
\{(0,0),(1,1),(2,1)\},
\]
the cross-support pair $(1,1)\in S_1$ and $(2,1)\in S_2$ has
\[
L(1,1)=L(2,1)=1,
\qquad
Q(1,1)=1\neq 2=Q(2,1).
\]
The compatibility condition \eqref{eq:cross-support-compatibility} fails, so $L$ does not realize $Q$ on $S_1\cup S_2$.
\end{example}

For cross-context mechanistic analysis, the two checks are therefore separate. Internal overlap should be computed at the cut where the localization is stated, since source overlap can be too small. A localization that is valid separately on two supports extends to their union exactly when the cross-support access collisions satisfy \eqref{eq:cross-support-compatibility}.
